% Part: computability
% Chapter: computability-theory
% Section: k-1

\documentclass[../../../include/open-logic-section]{subfiles}

\begin{document}

\olfileid[hi]{cmp}{thy}{k1}
\olsection{अपचयनीयता का एक उदाहरण}

आइए \olref[ppr]{prop:reduce} के एक अनुप्रयोग पर विचार करें।

\begin{prop}
\ollabel{prop:k1}
मानें
\[
K_1 = \Setabs{e}{\cfind{e}(0) \fdefined}.
\]
तब $K_1$ संगणनीयतः परिगणनीय है, किंतु संगणनीय नहीं।
\end{prop}

\begin{proof}
चूँकि $K_1 = \Setabs{e}{\lexists[s][T(e,0,s)]}$, इसलिए
\olref[eqc]{thm:exists-char} से $K_1$ संगणनीयतः परिगणनीय है।

$K_1$ को असंगणनीय दिखाने के लिए हम दिखाएँ कि $K_0$ उसमें अपचयनीय है।

\begin{explain}
यह कुछ कठिन है, क्योंकि $K_1$ का उपयोग करके हम केवल उन संगणनाओं के
विषय में प्रश्न पूछ सकते हैं जो किसी विशेष आगत $0$ से आरंभ होती हैं।
मान लें कि आपका कोई बुद्धिमान मित्र ऐसे प्रश्नों का उत्तर दे सकता है
(ऐसे मित्रों को ``दैवज्ञ'' कहा जाता है)। अब कोई आपसे पूछता है कि
$\tuple{e,
  x}$, $K_0$ में है या नहीं, अर्थात् मशीन $e$, आगत $x$ पर
रुकती है या नहीं। आप एक दूसरी मशीन $e_x$ बना सकते हैं, जो \emph{किसी
भी} आगत को अनदेखा करके~$e$ को आगत $x$ पर चलाती है। स्पष्टतः मशीन $e$ का
आगत $x$ पर रुकना इस प्रश्न के समतुल्य है कि मशीन $e_x$ आगत $0$ (या किसी
अन्य आगत) पर रुकती है या नहीं। अतः आप अपने मित्र से पूछते हैं कि यह नई
मशीन $e_x$, आगत $0$ पर रुकती है या नहीं; परिवर्तित प्रश्न का उसका उत्तर
मूल प्रश्न का उत्तर देता है। इससे $K_0$ से~$K_1$ में इच्छित अपचयन मिलता है।
\end{explain}

सार्वत्रिक आंशिक संगणनीय फलन का उपयोग करके $f$ को इस प्रकार परिभाषित
3-स्थानी फलन मानें:
\[
f(x,y,z) \simeq \cfind{x}(y).
\]
ध्यान दें कि $f$ अपनी तीसरी आगत को पूरी तरह अनदेखा करता है। ऐसा सूचकांक
$e$ चुनें कि $f = \cfind{e}[3]$; अतः हमें
\[
\cfind{e}[3](x,y,z) \simeq \cfind{x}(y).
\]
$s$-$m$-$n$ प्रमेय से कोई फलन $s(e,x,y)$ ऐसा है कि प्रत्येक $z$ के लिए
\begin{align*}
\cfind{s(e,x,y)}(z) & \simeq \cfind{e}[3](x,y,z) \\
& \simeq \cfind{x}(y).
\end{align*}

\begin{explain}
ऊपर के अनौपचारिक तर्क के पदों में, $s(e,x,y)$ उस मशीन का सूचकांक है जो
किसी भी आगत $z$ को अनदेखा करके $\cfind{x}(y)$ संगणित करती है।
\end{explain}

विशेषतः, हमें
\[
\cfind{s(e,x,y)}(0) \fdefined \quad \text{तभी और केवल तभी} \quad
\cfind{x}(y) \fdefined.
\]
दूसरे शब्दों में, $\tuple{x, y} \in K_0$ तभी और केवल तभी जब
$s(e,x,y) \in
K_1$। अतः इस प्रकार परिभाषित फलन $g$
\[
g(w) = s(e,(w)_0,(w)_1)
\]
$K_0$ से~$K_1$ में एक अपचयन है।
\end{proof}

\end{document}
